\documentclass[conference]{IEEEtran}
\IEEEoverridecommandlockouts

\usepackage{amsfonts}
\usepackage{amssymb}
\usepackage{amsmath}
\usepackage{booktabs}
\usepackage{multirow}
\usepackage{float}
\usepackage{tabularx}
\usepackage{graphicx}
\usepackage{xurl}
\usepackage[hidelinks]{hyperref}

\def\BibTeX{{\rm B\kern-.05em{\sc i\kern-.025em b}\kern-.08em
    T\kern-.1667em\lower.7ex\hbox{E}\kern-.125emX}}

\begin{document}

\title{Secure Digital Voice Intercom System Using Python-Based Audio Streaming}

\author{Minh~D.~H.~Bui\IEEEauthorrefmark{1},
        Khang~T.~Phan\IEEEauthorrefmark{1},
        Linh~N.~Do\IEEEauthorrefmark{1},
        Thien~D.~Do\IEEEauthorrefmark{1},
        Thanh~C.~Nguyen\IEEEauthorrefmark{1}\\
        \IEEEauthorblockA{\IEEEauthorrefmark{1}School of Electrical and Electronic Engineering, Hanoi University of Science and Technology, Hanoi 100000, Vietnam.}
}

\maketitle

\begin{abstract}
This paper presents a secure digital voice intercom system implemented with a Python-based audio relay and browser clients. The target deployment is a local-area network in which client machines should join the intercom without installing Python source code or native executables. The server hosts an HTTPS web interface and a WebSocket Secure endpoint, while each browser client captures microphone audio, converts it to mono 16-kHz PCM16 frames, and transmits binary audio frames to the server. The server admits clients through a room-key check and relays audio only to other clients in the same logical room. The paper develops the audio representation, relay fan-out, access-control, latency, concurrency, and measurement models that correspond to the implemented system. It also defines a reproducible experimental procedure for validating audio flow, relay scalability, room isolation, authentication rejection, and playback stability. The resulting design demonstrates a practical browser-accessible intercom architecture with low client deployment overhead and measurable real-time behavior.
\end{abstract}

\begin{IEEEkeywords}
Digital intercom, Python, WebSocket Secure, browser audio, real-time communication, LAN voice communication, secure audio streaming.
\end{IEEEkeywords}

\section{Introduction}

Internal voice communication remains necessary in laboratories, classrooms, buildings, factories, and control rooms where short spoken messages must be exchanged with low delay. Conventional analog intercoms are simple and robust, but they are difficult to scale, integrate with software, or instrument for measurement. IP-based intercom systems address these limitations by moving capture, access control, transport, and monitoring into software running over existing network infrastructure.

A digital intercom must satisfy two requirements simultaneously. First, the audio path must be sufficiently responsive for natural conversation. Previous studies on VoIP and real-time conferencing show that delay, jitter, loss, and wireless access conditions strongly affect user experience \cite{ieee_voip_delay_jitter,ieee_packet_loss_qoe,ieee_surecall}. Second, the system must prevent unauthorized users from joining the communication room and must protect the audio channel from passive observation. Security work on WebRTC signaling shows that real-time communication systems require authenticated session establishment to prevent impersonation and man-in-the-middle attacks \cite{ieee_webrtc_signaling_security}.

The system in this paper addresses a specific but practical deployment problem: client machines should be able to join an intercom room without receiving project files or installing a native client. To achieve this, the proposed design uses a Python server and browser clients. The Python server provides an HTTPS page and a WebSocket Secure (WSS) relay. The browser performs microphone capture and audio packetization, and the server performs room-key admission, room membership management, and audio forwarding.

The contributions of this paper are:
\begin{enumerate}
    \item a browser-accessible LAN intercom architecture that removes native client distribution;
    \item an audio transport model based on browser capture, 16-kHz PCM16 packetization, WSS transport, and Python room relay;
    \item a security and access-control model based on HTTPS/WSS transport and room-key admission;
    \item a measurement and experimental framework for evaluating connectivity, traffic load, relay scalability, authentication behavior, and playback stability.
\end{enumerate}

The paper is organized into five sections. Section II presents the methodology. Section III defines the system model. Section IV describes the experimental procedure. Section V concludes the paper.

\section{Methodology}

\subsection{Architectural Rationale}

The system is designed for a trusted LAN environment. One machine acts as the intercom server, while all participants connect through browsers. This architectural choice follows directly from the deployment requirement: if clients should not receive source code or executables, then the browser is the only realistic runtime already present on most client machines. The server is therefore responsible for hosting the interface, enforcing admission control, maintaining room state, and relaying media frames.

\begin{figure}[t]
    \centering
    \fbox{\begin{minipage}[c][1.45in][c]{0.92\columnwidth}
    \centering
    Placeholder for Fig. 1.\\
    Required content: browser clients, HTTPS page delivery, WSS connection, room-key admission, room-based audio relay.
    \end{minipage}}
    \caption{Architecture of the proposed browser-based secure digital intercom. The figure is meaningful because it maps each implementation component to one communication role.}
    \label{fig:architecture}
\end{figure}

Fig.~\ref{fig:architecture} should show only components that exist in the implementation: browser clients, HTTPS service, WSS endpoint, room-key validation, room table, and audio relay. It should not include RTP, LDPC, native audio libraries, or peer-to-peer media paths because those are not part of the implemented project.

\subsection{Audio Representation}

The browser microphone produces a discrete audio stream at the browser audio-context rate $f_c$, usually 44.1 kHz or 48 kHz. The implementation standardizes the outgoing stream to
\begin{equation}
    f_s = 16\,\mathrm{kHz},
    \label{eq:fs}
\end{equation}
because 16-kHz mono speech provides adequate intelligibility for an intercom while keeping bandwidth lower than full-bandwidth browser audio.

For each captured block, let $x[n]\in[-1,1]$ be the normalized browser audio samples. The client maps the block to output samples $x_s[m]$ by interval averaging:
\begin{equation}
    x_s[m] = \frac{1}{|\mathcal{I}_m|}
    \sum_{n \in \mathcal{I}_m} x[n],
    \label{eq:resampling}
\end{equation}
where $\mathcal{I}_m$ is the set of input indices assigned to output sample $m$. This equation is included because it describes the actual downsampling strategy in the browser client rather than an abstract codec not implemented by the system.

The resampled sequence is quantized to signed 16-bit PCM:
\begin{equation}
    q[m]=
    \begin{cases}
        \lfloor 32767x_s[m]\rfloor, & x_s[m]\geq 0,\\
        \lceil 32768x_s[m]\rceil, & x_s[m]<0.
    \end{cases}
    \label{eq:pcm}
\end{equation}
Thus, the nominal payload rate per continuously speaking client is
\begin{equation}
    R_{\mathrm{pcm}} = f_s b_q = 16000 \times 16 = 256\,\mathrm{kbps},
    \label{eq:payloadrate}
\end{equation}
where $b_q=16$ bits/sample. This rate is useful because it gives a baseline against which measured WebSocket/TLS traffic can be interpreted.

\subsection{Relay Method}

WebSocket is used because it provides a persistent full-duplex browser--server channel and has been evaluated for real-time communication scenarios \cite{ieee_websocket_realtime}. In this project, control messages are JSON text frames and audio messages are binary PCM frames. The relay method is room-based. Let $C_r(t)$ be the set of clients in room $r$ at time $t$. When client $c_i$ sends frame $a_i(k)$, the forwarding set is
\begin{equation}
    \mathcal{R}_i(t) = C_r(t)\setminus\{c_i\}.
    \label{eq:relayset}
\end{equation}
Equation~(\ref{eq:relayset}) is meaningful because it defines both room isolation and echo avoidance: the sender does not receive its own frame from the server, and clients in other rooms are excluded.

If $N=|C_r(t)|$ clients are active in one room and all transmit at approximately $R_{\mathrm{pcm}}$, the server receives
\begin{equation}
    R_{\mathrm{in}} = NR_{\mathrm{pcm}},
    \label{eq:rin}
\end{equation}
and relays
\begin{equation}
    R_{\mathrm{out}} = N(N-1)R_{\mathrm{pcm}}.
    \label{eq:rout}
\end{equation}
Equations~(\ref{eq:rin}) and~(\ref{eq:rout}) motivate the multi-client experiment because they predict that relay bandwidth grows quadratically with room size.

\subsection{Security Model}

The security model is intentionally scoped to a controlled LAN. Transport protection is provided by HTTPS/WSS, and admission control is enforced through a shared room key. Let $K_s$ be the configured server key and $K_c$ be the key submitted by a client. The admission function is
\begin{equation}
    A(c)=
    \begin{cases}
        1, & K_c=K_s,\\
        0, & K_c\neq K_s.
    \end{cases}
    \label{eq:admission}
\end{equation}
Only clients with $A(c)=1$ are inserted into the server room state. This formula is not a cryptographic primitive; it formalizes the implemented access-control rule and provides a direct basis for the authentication-rejection experiment.

The server is trusted and can access relayed audio payloads. Therefore, the system provides encrypted browser--server transport and access control, but not end-to-end confidentiality against the relay server. More advanced peer-to-peer designs require stronger signaling protection and key management, as discussed in WebRTC security literature \cite{ieee_webrtc_signaling_security}.

\subsection{Latency and Playback Stability}

The one-way delay of a browser audio frame is decomposed as
\begin{equation}
    D_{\mathrm{total}} =
    D_{\mathrm{cap}} + D_{\mathrm{conv}} + D_{\mathrm{net}}
    + D_{\mathrm{relay}} + D_{\mathrm{play}},
    \label{eq:latency}
\end{equation}
where $D_{\mathrm{cap}}$ is browser capture callback delay, $D_{\mathrm{conv}}$ is sample-rate conversion and PCM packing time, $D_{\mathrm{net}}$ is WSS/TCP/TLS transport delay, $D_{\mathrm{relay}}$ is server scheduling and forwarding delay, and $D_{\mathrm{play}}$ is receiver-side playback scheduling delay. This model is used to reason about where delay can enter the implemented system.

Because the browser client uses WebSocket over TCP, packet loss is not directly exposed like it would be in a UDP/RTP implementation. Congestion or scheduling delay is instead reflected in queueing and playback delay. The browser therefore estimates the playback queue as
\begin{equation}
    Q_{\mathrm{play}}(t) =
    \max(0,T_{\mathrm{next}}-T_{\mathrm{ctx}}),
    \label{eq:qplay}
\end{equation}
where $T_{\mathrm{next}}$ is the next scheduled playback time and $T_{\mathrm{ctx}}$ is the current audio-context time. A bounded $Q_{\mathrm{play}}(t)$ indicates stable playback; monotonic growth indicates accumulating delay.

\section{System Model}

\subsection{Server State and Events}

The server state is the set
\begin{equation}
    S(t)=\{(id_i,name_i,room_i,ws_i,\tau_i)\}_{i=1}^{M(t)},
    \label{eq:serverstate}
\end{equation}
where $id_i$ is a generated client identifier, $name_i$ is a display name, $room_i$ is the room, $ws_i$ is the WebSocket connection, and $\tau_i$ is the join time. This state model is included because every server decision is derived from it: presence updates, room filtering, disconnect handling, and forwarding.

The server processes four event types:
\begin{enumerate}
    \item \textbf{Page request}: return the browser interface through HTTPS.
    \item \textbf{Join message}: validate the room key using (\ref{eq:admission}) and insert the client into $S(t)$ if accepted.
    \item \textbf{Audio frame}: validate payload size and relay the frame according to (\ref{eq:relayset}).
    \item \textbf{Disconnect}: remove the client from $S(t)$ and update room presence.
\end{enumerate}

\begin{figure}[t]
    \centering
    \fbox{\begin{minipage}[c][1.45in][c]{0.92\columnwidth}
    \centering
    Placeholder for Fig. 2.\\
    Required content: join validation, room insertion, binary audio frame, relay to room peers, disconnect removal.
    \end{minipage}}
    \caption{Server event model. The figure should correspond to the four server event types described in Section III-A.}
    \label{fig:serverevents}
\end{figure}

\subsection{Browser Client Pipeline}

The browser client pipeline has four stages: capture, conversion, transmission, and playback. Capture is initiated only after the browser grants microphone permission. Conversion applies (\ref{eq:resampling}) and (\ref{eq:pcm}). Transmission sends PCM frames as WSS binary frames. Playback converts received PCM16 frames back to floating-point samples and schedules them on the local audio context.

The client also maintains local counters for captured frames, transmitted bytes, received bytes, played packets, capture errors, and playback queue duration. These counters are necessary because the server cannot infer browser capture failures or local playback delay from network traffic alone.

\subsection{Concurrency and Capacity}

The Python server uses asynchronous event handling. Let $\lambda_a$ be the incoming audio-frame rate and $\mu_r$ be the relay service rate. A necessary stability condition is
\begin{equation}
    \rho = \frac{\lambda_a}{\mu_r} < 1.
    \label{eq:rho}
\end{equation}
If $\rho$ approaches or exceeds unity, frames will accumulate in buffers, increasing $D_{\mathrm{relay}}$ in (\ref{eq:latency}) and potentially increasing $Q_{\mathrm{play}}(t)$ in (\ref{eq:qplay}). Thus, relay throughput and playback queue duration must be evaluated together.

\subsection{Measurement Variables}

Table~\ref{tab:variables} contains only variables used by the experimental procedure. The table is included to make the evaluation reproducible and to prevent ambiguity between server-side and browser-side observations.

\begin{table}[t]
\centering
\caption{Measurement variables used in the experimental procedure}
\label{tab:variables}
\renewcommand{\arraystretch}{1.15}
\begin{tabularx}{\columnwidth}{lX}
\toprule
\textbf{Variable} & \textbf{Definition and Purpose} \\
\midrule
$N$ & Number of accepted browser clients in one room; used for relay scalability. \\
$B_{\mathrm{rx}}$ & Server-side received audio bitrate; compared with (\ref{eq:rin}). \\
$B_{\mathrm{relay}}$ & Server-side relayed audio bitrate; compared with (\ref{eq:rout}). \\
$P_{\mathrm{rx}}$ & Server-side received audio packet rate; indicates audio frame arrival continuity. \\
$P_{\mathrm{relay}}$ & Server-side relayed packet rate; indicates whether frames are forwarded. \\
$F_{\mathrm{auth}}$ & Failed room-key attempts; validates access-control behavior. \\
$F_{\mathrm{drop}}$ & Server-rejected audio frames; detects invalid or oversized payloads. \\
$Q_{\mathrm{play}}$ & Browser-side playback queue; evaluates receiver stability. \\
$E_{\mathrm{cap}}$ & Browser-side capture errors; detects microphone permission/device failures. \\
\bottomrule
\end{tabularx}
\end{table}

\section{Experimental Procedure}

\subsection{Testbed and Configuration}

The experiment uses one Python server and at least two browser clients on the same LAN. The server binds to TCP port 8443 and serves the interface through HTTPS. Each client uses a modern browser, grants microphone permission, enters a display name, selects a room, and submits the shared room key. Headphones are used to reduce acoustic echo during full-duplex tests.

\begin{figure}[t]
    \centering
    \fbox{\begin{minipage}[c][1.45in][c]{0.92\columnwidth}
    \centering
    Placeholder for Fig. 3.\\
    Required content: one Python server, multiple browser clients, LAN/Wi-Fi links, and measured server/client variables.
    \end{minipage}}
    \caption{LAN testbed for evaluating the proposed intercom. The figure supports the experimental procedure by identifying where measurements are collected.}
    \label{fig:testbed}
\end{figure}

The server is started with a fixed reporting interval:
\begin{verbatim}
python -m web_intercom.server --host 0.0.0.0
  --port 8443 --stats-interval 5
\end{verbatim}
The same configuration is used across scenarios so that differences in measurements arise from room size, room allocation, and authentication behavior rather than from server reconfiguration.

\subsection{Scenarios}

Four scenarios are used because each one validates a distinct system property:
\begin{enumerate}
    \item \textbf{Two-client audio flow}: two clients join the same room and exchange speech. This validates the minimum useful intercom case.
    \item \textbf{Multi-client relay}: $N\geq3$ clients join one room. This validates whether measured relay traffic follows the fan-out model in (\ref{eq:rout}).
    \item \textbf{Room isolation}: clients are split into two rooms. This validates that (\ref{eq:relayset}) prevents cross-room audio delivery.
    \item \textbf{Authentication rejection}: a client submits an incorrect room key. This validates that (\ref{eq:admission}) rejects unauthorized access and increments $F_{\mathrm{auth}}$.
\end{enumerate}

\subsection{Data Processing}

For each scenario, measurements are collected over duration $T$. The mean received bitrate is
\begin{equation}
    \bar{B}_{\mathrm{rx}} =
    \frac{8\,bytes_{\mathrm{rx}}(T)}{1000T},
    \label{eq:meanrx}
\end{equation}
and the peak interval received bitrate is
\begin{equation}
    B_{\mathrm{rx,peak}} =
    \max_k B_{\mathrm{rx}}(\Delta t_k).
    \label{eq:peakrx}
\end{equation}
The relay expansion factor is
\begin{equation}
    \Gamma =
    \frac{bytes_{\mathrm{relay}}(T)}
         {bytes_{\mathrm{rx}}(T)}.
    \label{eq:gamma}
\end{equation}
In a single room where each client transmits comparable audio, $\Gamma$ should approach $N-1$. A much smaller value indicates that frames are not being forwarded to all intended peers; a much larger value would indicate duplicate forwarding.

Playback stability is evaluated from $Q_{\mathrm{play}}(t)$. A run is stable when $Q_{\mathrm{play}}(t)$ remains bounded during the observation window. This criterion is consistent with prior work showing that delay and jitter are central to packet voice quality \cite{ieee_voip_delay_jitter,ieee_jitter_buffer_pd}.

\subsection{Acceptance Criteria}

A run is considered valid when the following conditions are satisfied:
\begin{enumerate}
    \item accepted clients appear in the intended room;
    \item browser sent-packet counters are non-zero during speech;
    \item server received and relayed packet counters are non-zero;
    \item browser received and played-packet counters are non-zero;
    \item room isolation prevents audio delivery across different rooms;
    \item incorrect room keys are rejected and counted by $F_{\mathrm{auth}}$;
    \item $Q_{\mathrm{play}}(t)$ remains bounded;
    \item $E_{\mathrm{cap}}$ and $F_{\mathrm{drop}}$ remain zero in normal operation.
\end{enumerate}
These criteria directly correspond to the proposed architecture: access control, room membership, relay forwarding, endpoint capture, endpoint playback, and stability.

\section{Conclusion}

This paper presented a secure digital voice intercom system using a Python-based audio relay and browser clients. The design is motivated by the requirement that client machines should connect without receiving source code or native executables. The methodology defined the implemented audio representation, relay rule, access-control function, latency decomposition, and playback-stability metric. The system model specified the server state, browser pipeline, asynchronous relay capacity condition, and measurement variables. The experimental procedure then mapped these models to concrete scenarios for validating audio flow, multi-client scalability, room isolation, authentication rejection, and playback stability.

The main trade-off is that browser/WSS deployment is easier than native UDP media transport, but it provides less direct control over packet timing and congestion behavior. The current system is therefore appropriate for LAN demonstrations and controlled experiments. Future work should investigate lower-latency browser processing, explicit one-way delay estimation, stronger end-to-end media protection, and comparison with a native UDP/RTP implementation under identical LAN conditions.

\bibliographystyle{IEEEtran}
\bibliography{ref}

\end{document}
